\newpage
\section{Proofs}\label{sec:tech:proof}

We omit the completely standard proof that basic typing
\(\basicvdash e : s\) is sound, assuming that all terms and qualifiers
in our typing rules and theorems are type-safe. Before presenting the
proof of the fundamental theorem and type soundness, we introduce
several useful lemmas.

\subsection{Lemmas}

\subsubsection{Common symbolic LTL$_f$ formulas}

\begin{lemma}\label{lemma:allTL}. $\allTL$ contains
  all well-formed traces.
  $\denot{\allTL} = \{ tr ~|~ \wellfoundedvdash tr \}$.
\end{lemma}

\begin{lemma}\label{lemma:espTL}. $\epsTL$ only contains the empty trace. $\denot{\epsTL} = \{ [] \}$.
\end{lemma}

\begin{lemma}\label{lemma:empTL} $\empTL$ contains no
  traces. $\denot{\empTL} = \emptyset$.
\end{lemma}

\subsubsection{Denotations}

\begin{lemma}\label{lemma:lastA-denot}[Denotation of singleton modality] For all symbolic event $\msgB{op}{\overline{x_i}}{\phi}$ and values $\overline{v_i}$,
{\small\begin{align*}
    \phi\msubst{x_i}{v_i} \implies
    [\zevent{op}{\overline{v_i}}] \in \denot{\lastA\msgB{op}{\overline{x_i}}{\phi}}
\end{align*}}
\end{lemma}

\begin{lemma}\label{lemma:seqA-denot}[Denotation of concatenation] For all automata $A_1$ and $A_2$ and trace $\alpha$,
{\small\begin{align*}
    \alpha \in \denot{A_1 \seqA A_2} \iff (\exists \alpha_1\;\alpha_2. \alpha = \alpha_1 \listconcat \alpha_2 \land \alpha_1 \in \denot{A_1} \land \alpha_2 \in \denot{A_2})
\end{align*}}
\end{lemma}

\begin{lemma}\label{lemma:choice-denot}[Denotation of choice] For all term $e_1$ and $e_2$ and \PAT $\tau$,
{\small\begin{align*}
    e_1 \in \denot{\tau} \land e_2 \in \denot{\tau} \impl e_1 \oplus e_2 \in \denot{\tau}
\end{align*}}
\end{lemma}

\begin{lemma}\label{lemma:pure-denot}[Denotation of pure computation] For all term $e_1$ and $e_2$ and \PAT $\tau$,
{\small\begin{align*}
    (\forall \alpha\;\beta. \msteptr{\alpha}{(\beta, e)}{[]}{(\beta, e')}) \impl
   e\in \denot{\tau} \iff e' \in \denot{\tau}
\end{align*}}
\end{lemma}

\begin{lemma}[Buffer Partition]\label{lemma:buf-partition} For all capability $\Theta$, automata $F$ and buffer $\beta$, we have
\begin{align*}
    \beta \in \denot{\Theta_1 \cup \Theta_2} \iff
    \exists \beta_1\;\beta_2. \beta_1 \cup \beta_2 = \beta \land \beta_1 \cap \beta_2 = \emptyset \land \beta_1 \in \Theta_1 \land \beta_2 \in \Theta_2
\end{align*}
\end{lemma}

% \begin{lemma}[Buffer Partition]\label{lemma:buf-partition} For all capability $\Theta$, automata $F$ and buffer $\beta$, we have
% \begin{align*}
%     (\Theta_1 \cup \Theta_2, \beta) \sim F \impl
%     \exists \beta_1\;\beta_2. \beta_1 \cup \beta_2 = \beta \land \beta_1 \cap \beta_2 = \emptyset \land (\Theta_1, \beta_1) \sim F \land (\Theta_2, \beta_2) \sim F
% \end{align*}
% \end{lemma}

% [\zevent{op}{\overline{\sigma(v_i)}}] \in \denot{\lastA\msg{op}{\sigma(\phi)\msubst{x_i}{v_i}}} 

\subsubsection{Subtyping}

\begin{lemma}\label{lemma:pure-subtyping}[Pure Subtyping Soundness] For Given type context $\Gamma$ and well-formed pure refinement type $t$ and $t'$:
  $\Gamma \vdash t <: t' \implies \forall \sigma \in \denot{\Gamma}. \denot{\sigma(t)} \subseteq \denot{\sigma(t')}$
\end{lemma}

\begin{lemma}\label{lemma:subtyping}[Subtyping Soundness] For Given type context $\Gamma$ and well-formed \PAT $\tau$ and $\tau'$:
  $\Gamma \vdash \tau <: \tau' \implies \forall \sigma \in \denot{\Gamma}. \denot{\sigma(\tau)} \subseteq \denot{\sigma(\tau')}$
\end{lemma}

\subsubsection{Substitution}

\begin{lemma}[Substitution Lemma]\label{lemma:subt} For Given type context $\Gamma$, variable $x$, well-formed pure refinement type $t$,\PAT $\tau$ and term $e$:
  $\Gamma, x{:}t;\Delta;\Theta \vdash e : \tau \impl \forall v. \Gamma \vdash v : t \impl \Gamma;\Delta;\Theta \vdash e[x\mapsto v] : \tau[x\mapsto v]$
\end{lemma}

\subsubsection{Handler Contexts}

\begin{definition}[Well-formed handler context]
  \label{lemma:built-in-typing-ap}
  The handler specification $\Delta$ is well-formed iff for all operator $\eff{op}$ and its \PAT $\overline{y{:}b} \garr \overline{x{:}t} \sarr \rg{H}{\lastA\msgB{op}{\overline{y}}{\phi}}{F}$ and capability $\{\overline{\eff{op_i}}\}$ in $\Delta$ satisfying
{\small\begin{align*}
    &\forall \overline{y{:}b}.\forall \alpha_h \in \denot{H}. \forall \overline{c \in \denotation{t}}. \forall \overline{c_{ij}}.\forall \overline{\alpha_i}. \alpha_1 \listconcat [\zevent{op_1}{\overline{c_{1j}}}] \listconcat ... [\zevent{op_n}{\overline{c_{nj}}}] \listconcat \alpha_{n+1} \in \denot{F} \impl
    \\&\qquad \alpha_h \vDash \zevent{op}{\overline{c}} \Downarrow \{\zevent{op_i}{\overline{c_{ij}}}\} \land \phi\msubst{x}{c}
\end{align*}}
\end{definition}

\begin{lemma}[Well-formed handler context with subsumption]
  \label{lemma:built-in-typing-sub}
For given well-formed handler specification $\Delta$, type context $\Gamma$, and effect operator $\eff{op}$
{\small\begin{align*}
    &\Delta(\eff{op}) = \langle \tau, \Theta \rangle \impl \Gamma \vdash \tau <: \overline{x{:}t} \sarr \rg{H}{\lastA\msgB{op}{\overline{y}}{\phi}}{F} \impl
    \\&\forall \sigma \in \denot{\Gamma}. \forall \alpha_h \in \denot{\sigma(H)}. \forall \overline{c \in \denotation{\sigma(t)}}. \forall \alpha_f \in \denot{\sigma(F)}.
    \exists \overline{\alpha_i}. \exists \overline{m_i}.\;
    \\&\alpha_1 \listconcat [m_1] \listconcat ... [m_{n}] \listconcat \alpha_{n+1} = \alpha_f  \land (\forall \eff{op_i}. \eff{op_i} \in \Theta {\iff} \exists \overline{c_i}.m_i = \zevent{op_i}{\overline{c_i}}) \impl \alpha_h \vDash \zevent{op}{\overline{c}} \Downarrow \{m_i\} \land \sigma(\phi)\msubst{x}{c}
\end{align*}}
\end{lemma}

\begin{lemma}[Well-formed pure context with subsumption]
  \label{lemma:built-in-typing-sub-pure}
For given specification $\Delta$, type context $\Gamma$, and pure operator $\primop$
{\small\begin{align*}
    &\Delta(\primop) = t \land \Gamma \vdash t <: \overline{y{:}t_y}\sarr t_x \impl \forall \sigma \in \denot{\Gamma}. \forall \overline{c_y \in \denot{\sigma(t_y)}}. \primop (\overline{c_y}) \Downarrow c \impl c \in \denot{\sigma(t_x\msubst{y}{c_y})}
\end{align*}}
\end{lemma}

\subsection{Fundamental Theorem}

We first prove the fundamental theorem for values.

\begin{theorem}\label{theorem:pure-fundamental-ap}[Pure Fundamental Theorem] For Given type context $\Gamma$ and well-formed value $v$ as well as pure refinement type $t$:
  $\Gamma \vdash v: t \implies \forall \sigma \in \denot{\Gamma}. \sigma(v) \in \denot{\sigma(t)}$
\end{theorem}
\begin{proof} We proceed by induction over our type judgment $\Gamma;\Delta;\Theta \vdash e : \tau$, which has two cases proved as following:
\begin{enumerate}[label=Case]
\item:
{\footnotesize
\mprooftr{23mm}
{$\eraserf{\Gamma} \basicvdash v : b$}
{TVal}
{$\Gamma \vdash v : \rawnuot{b}{\vnu = v}$}
\ \\ \ \\}
where we need to prove $\forall \sigma \in \denot{\Gamma}. \sigma(v) \in\denot{\rawnuot{b}{\vnu = \sigma(v)}}$, which can be directly proved by definition of type denotation.
\item:
{\footnotesize
\mprooftr{38mm}
{$\Gamma \vdash v : t$ \quad
$\Gamma \vdash t <: t'$}
{TPureSub}
{$\Gamma \vdash v : t'$}
\ \\ \ \\}
where we have inductive hypothesis $\forall \sigma \in \denot{\Gamma}. \sigma(v) \in\denot{\sigma(t)}$ and need to prove $\forall \sigma \in \denot{\Gamma}. \sigma(v) \in\denot{\sigma(t')}$, which can be directly proved by soundness lemma of pure subtyping (lemma~\ref{lemma:pure-subtyping}).
\end{enumerate}
\end{proof}


The fundamental theorem for a controller program consists of two parts: (1) the history, current, and future traces of a well-typed term \( e \) are consistent with the corresponding \PAT; (2) the realizability guarantee provided by the capability. We first prove the first part, as follows.

\begin{theorem}\label{theorem:fundamental-ap}[Fundamental Theorem For Trace Consistency] Given
  a well-formed handler specification $\Delta$,
    the trace of effects produced by a well-typed term
    $e$ is captured by its corresponding \PAT $\tau$:
  $\Gamma; \Delta; \Theta \vdash e : \tau \implies \forall \sigma, \sigma \in
    \denotation{\Gamma} \impl \sigma(e) \in
    \denotation{\sigma(\tau)}$.
\end{theorem}
\begin{proof} We proceed by induction over our type judgment $\Gamma;\Delta;\Theta \vdash e : \tau$, which has $8$ cases proved as following:
\begin{enumerate}[label=Case]

\item:
{\footnotesize
\mprooftr{87mm}
{
$\Delta(\eff{op}) =\langle \mykw{gen}\;\tau, \Theta' \rangle$ \quad
$\Gamma \vdash \tau <: \overline{x_i{:}t_i}\sarr
\rg{H}{\lastA\msg{op}{\phi}}{A \seqA F}$ \quad
$\forall i. \Gamma \vdash v_i : t_i$ \quad
 $\Gamma; \Delta; \Theta \cup \Theta' \vdash e : \rg{H \seqA \lastA\msg{op}{\phi\msubst{x_i}{v_i}}}{A}{F}$
 }
{TGen}
{$\Gamma; \Delta; \Theta \vdash \zgen{op}{\overline{v_i}}{e} : \rg{H}{\lastA\msg{op}{\phi\msubst{x_i}{v_i}} \seqA A}{F}$}
\ \\ \ \\}
This rule assume that $e \equiv \zgen{op}{\overline{v}}{e}, \tau \equiv \rg{H}{\lastA\msg{op}{\phi\msubst{x}{v}} \seqA A}{F}$, thus we need to prove
{\footnotesize\begin{align*}
    \forall \sigma, \sigma \in \denot{\Gamma} \impl \sigma(\zgen{op}{\overline{v}}{e}) \in \denot{\sigma(\rg{H}{\lastA\msg{op}{\phi\msubst{x}{v}} \seqA A}{F})}
\end{align*}}\noindent
From the induction hypothesis and the precondition of this rule, we have
{\footnotesize\begin{alignat}{2}
\setcounter{equation}{0}
&\Delta(\eff{op}) =\langle \mykw{gen}\;\tau, \Theta' \rangle \ptag{assumption} \label{found:1-1} \\
&\Gamma \vdash \tau <: \overline{x{:}t}\sarr \rg{H}{\lastA\msg{op}{\phi}}{A \seqA F} \ptag{assumption} \label{found:1-2} \\
&\forall i. \Gamma \vdash v_i : t_i \ptag{assumption} \label{found:1-3} \\
% &\Gamma \vdash \nuot{unit}{\top} <: \nuot{unit}{\phi\msubst{x_i}{v_i}} \ptag{assumption} \label{found:1-4} \\
&\Gamma ~|~ \Theta \cup \Theta' \vdash e : \rg{H \seqA \lastA\msg{op}{\phi\msubst{x}{v}}}{A}{F} \ptag{assumption} \label{found:1-5} \\
&\forall \sigma \in \denot{\Gamma}. \sigma(e) \in \denot{\sigma(\rg{H \seqA \lastA\msg{op}{\phi\msubst{x}{v}}}{A}{F})} \ptag{induction hypothesis} \label{found:1-6} \\
&\forall i. \forall \sigma \in \denot{\Gamma}. \sigma(v_i) \in  \sigma(\denot{t_i}) \ptag{\ref{found:1-3} and Lemma~\ref{theorem:pure-fundamental-ap}} \label{found:1-7} \\
&\forall \sigma \in \denot{\Gamma}. \sigma(\phi)\msubst{x_i}{v_i} \ptag{Lemma~\ref{lemma:built-in-typing-sub}, \ref{found:1-1}, \ref{found:1-2}, and \ref{found:1-3}} \label{found:1-8}
\end{alignat}}\noindent
According to denotation of \PAT and assumption~\ref{found:1-6}, we have
{\footnotesize\begin{alignat}{2}
&\forall \sigma \in \denot{\Gamma}. \forall \alpha_h \in\denot{\sigma(H \seqA \lastA\msg{op}{\phi\msubst{x}{v}})}.\forall \alpha_f \in\denot{F}. \forall \alpha\;\beta\;\beta'\;e_h\;e_f. \notag
\\&\quad \msteptr{[]}{(\emptyset, e_h)}{\alpha_h}{(\beta, ())} \land \msteptr{\alpha_h}{(\beta, \sigma(e))}{\alpha}{(\beta', ())} \land \msteptr{\alpha_h \listconcat \alpha }{(\beta', e_f)}{\alpha_f}{(\emptyset, ())} \impl \notag
\\&\quad \alpha \in\denot{\sigma(A)}) \ptag{assumption \ref{found:1-6}} \label{found:1-9}
\end{alignat}}\noindent
From now, we consider each $\sigma \in \denot{\Gamma}$, and try to prove the subgoal of this case, i.e.,
{\footnotesize\begin{align*}
    \sigma(\zgen{op}{\overline{v}}{e}) \in \denot{\sigma(\rg{H}{\lastA\msg{op}{\phi\msubst{x}{v}} \seqA A}{F})}
\end{align*}}\noindent
According to denotation of \PAT, we need to prove for all $\forall \alpha_h\;\alpha_f\;\alpha\;\beta\;\beta'\;e_h\;e_f.$ where $\alpha_h \in \denot{\sigma(H)}$ and $\alpha_f \in \denot{\sigma(F)}$,
{\footnotesize\begin{align*}
&\msteptr{[]}{(\emptyset, e_h)}{\alpha_h}{(\beta, ())} \land \msteptr{\alpha_h}{(\beta, \zgen{op}{\overline{\sigma(v_i)}}{\sigma(e)}))}{\alpha}{(\beta', ())} \land \msteptr{\alpha_h \listconcat \alpha }{(\beta', e_f)}{\alpha_f}{(\emptyset, ())} \impl \notag
\\&\quad \alpha \in\denot{\sigma(\lastA\msg{op}{\phi\msubst{x}{v}} \seqA A)}
\end{align*}}\noindent
Then we have
{\footnotesize\begin{alignat}{2}
&\sigma \in \denot{\Gamma} \land \alpha_h \in \denot{\sigma(H)} \land  \alpha_f \in \denot{\sigma(F)} \ptag{assumption} \label{found:1-10}
\\& \msteptr{[]}{(\emptyset, e_h)}{\alpha_h}{(\beta, ())} \ptag{assumption} \label{found:1-11}
\\& \msteptr{\alpha_h}{(\beta, \zgen{op}{\overline{\sigma(v_i)}}{\sigma(e)}))}{\alpha}{(\beta', ())} \ptag{assumption} \label{found:1-12}
\\& \msteptr{\alpha_h \listconcat \alpha }{(\beta', e_f)}{\alpha_f}{(\emptyset, ())} \ptag{assumption} \label{found:1-13}
\\& [\zevent{op}{\overline{\sigma(v_i)}}] \in \denot{\lastA\msg{op}{\sigma(\phi)\msubst{x_i}{v_i}}} \ptag{lemma~\ref{lemma:lastA-denot}} \label{found:1-14}
\\& \alpha_h \listconcat [\zevent{op}{\overline{\sigma(v_i)}}] \in \denot{\sigma(H \seqA \lastA\msg{op}{\phi\msubst{x_i}{v_i}})} \ptag{lemma~\ref{lemma:seqA-denot}, \ref{found:1-12}, and \ref{found:1-11}} \label{found:1-15}
\\& \exists \alpha'. \alpha = \zevent{op}{\overline{\sigma(v_i)}} \cons \alpha' \land
\alpha_h \models \zevent{op}{\overline{\sigma(v_i)}} \Downarrow \beta_{\eff{op}} \land \notag
\\&\quad\msteptr{\alpha_h \listconcat [\zevent{op}{\overline{\sigma(v_i)}}]}{(\beta \cup \beta_{\eff{op}}, \sigma(e))}{\alpha'}{(\beta', ())} \ptag{\textsc{StGen} and \ref{found:1-12}} \label{found:1-16}
\end{alignat}}\noindent
Now, we can apply hypothesis~\ref{found:1-9} with
{\footnotesize\begin{align*}
    &\sigma \mapsto \sigma \quad \alpha_h \mapsto \alpha_h \listconcat [\zevent{op}{\overline{\sigma(v_i)}}] \quad \alpha_f \mapsto \alpha_f \quad \alpha \mapsto \alpha' \quad \beta \mapsto \beta \cup \beta_{\eff{op}} \quad
    e_h \mapsto e_h; \zgen{op}{\overline{\sigma(v_i)}}{()} \quad e_f \mapsto e_f
\end{align*}}\noindent
Then we have
{\footnotesize\begin{alignat}{2}
&\alpha' \in\denot{\sigma(A)} \ptag{hypothesis~\ref{found:1-9} with \ref{found:1-10}, \ref{found:1-11}, \ref{found:1-13}, \ref{found:1-15}, \ref{found:1-16}} \label{found:1-17}
\\&[\zevent{op}{\overline{\sigma(v_i)}}] \listconcat \alpha' \in\denot{\sigma(\lastA\msg{op}{\phi\msubst{x}{v}} \seqA A)}\ptag{hypothesis~\ref{found:1-17}} \label{found:1-18}
\end{alignat}}\noindent
that is sufficient to prove subgoal of this case.

\item:
{\footnotesize
\mprooftr{80mm}
{
  $\Delta(\eff{op}) =\langle \mykw{obs}\;\tau, \Theta'\rangle$ \quad
$\Gamma \vdash \tau <: \overline{x_i{:}t_i}\sarr
\rg{H}{\lastA\msgB{op}{\overline{y}}{\phi \land \overline{y = x}}}{A \seqA F}$ \quad
 $\Gamma, \overline{x{:}t}; \Delta; \Theta \cup \Theta' \vdash e : \rg{H \seqA \lastA\msgB{op}{\overline{y}}{\phi \land \overline{y = x}}}{A}{F}$
 }
{TObs}
{$\Gamma; \Delta; \{\eff{op}\} \cup \Theta \vdash \zobs{\overline{x}}{op}{e} : \rg{H}{\lastA\msg{op}{\phi}\seqA A}{F}$}
\ \\ \ \\}
This rule assume that $e \equiv \zobs{\overline{x}}{op}{e}, \tau \equiv\rg{H}{\lastA\msg{op}{\phi} \seqA A}{F}$, thus we need to prove
{\footnotesize\begin{align*}
    \forall \sigma, \sigma \in \denot{\Gamma} \impl \sigma(\zobs{\overline{x}}{op}{e}) \in \denot{\sigma(\rg{H}{\lastA\msg{op}{\phi} \seqA A}{F})}
\end{align*}}\noindent
From the induction hypothesis and the precondition of this rule, we have
{\footnotesize\begin{alignat}{2}
\setcounter{equation}{0}
&\Delta(\eff{op}) =\langle \mykw{obs}\;\tau, \Theta' \rangle \ptag{assumption} \label{found:2-1} \\
&\Gamma \vdash \tau <: \overline{x_i{:}t_i}\sarr\rg{H}{\lastA\msgB{op}{\overline{y}}{\phi \land \overline{y = x}}}{A \seqA F} \ptag{assumption} \label{found:2-2} \\
&\Gamma, \overline{x{:}t}; \Delta; \Theta \cup \Theta' \vdash e : \rg{H \seqA \lastA\msgB{op}{\overline{y}}{\phi \land \overline{y = x}}}{A}{F} \ptag{assumption} \label{found:2-3} \\
&\forall \sigma \in \denot{\Gamma, \overline{x{:}t}}. \sigma(e) \in \denot{\sigma(\rg{H \seqA \lastA\msgB{op}{\overline{y}}{\phi \land \overline{y = x}}}{A}{F})} \ptag{induction hypothesis} \label{found:2-4}
\end{alignat}}\noindent
According to denotation of \PAT and assumption~\ref{found:2-4}, we have
{\footnotesize\begin{alignat}{2}
&\forall \sigma \in\denot{\Gamma, \overline{x{:}t}}. \forall \alpha_h \in\denot{\sigma(H \seqA \lastA\msgB{op}{\overline{y}}{\phi \land \overline{y = x}})}.\forall \alpha_f \in\denot{F}. \forall \alpha\;\beta\;\beta'\;e_h\;e_f. \notag
\\&\quad \msteptr{[]}{(\emptyset, e_h)}{\alpha_h}{(\beta, ())} \land \msteptr{\alpha_h}{(\beta, \sigma(e))}{\alpha}{(\beta', ())} \land \msteptr{\alpha_h \listconcat \alpha }{(\beta', e_f)}{\alpha_f}{(\emptyset, ())} \impl \notag
\\&\quad \alpha \in\denot{\sigma(A)}) \ptag{assumption \ref{found:2-4}} \label{found:2-5}
\end{alignat}}\noindent
From now, we consider each $\sigma \in \denot{\Gamma}$, and try to prove the subgoal of this case, i.e.,
{\footnotesize\begin{align*}
    \sigma(\zobs{\overline{x}}{op}{e}) \in \denot{\sigma(\rg{H}{\lastA\msg{op}{\phi} \seqA A}{F})}
\end{align*}}\noindent
According to denotation of \PAT, we need to prove for all $\forall \alpha_h\;\alpha_f\;\alpha\;\beta\;\beta'\;e_h\;e_f.$ where $\alpha_h \in \denot{\sigma(H)}$ and $\alpha_f \in \denot{\sigma(F)}$,
{\footnotesize\begin{align*}
&\msteptr{[]}{(\emptyset, e_h)}{\alpha_h}{(\beta, ())} \land \msteptr{\alpha_h}{(\beta, \zobs{\overline{x}}{op}{\sigma(e)}))}{\alpha}{(\beta', ())} \land \msteptr{\alpha_h \listconcat \alpha }{(\beta', e_f)}{\alpha_f}{(\emptyset, ())} \impl \notag
\\&\quad \alpha \in\denot{\sigma(\lastA\msg{op}{\phi} \seqA A)}
\end{align*}}\noindent
Then we have
{\footnotesize\begin{alignat}{2}
&\sigma \in \denot{\Gamma} \land \alpha_h \in \denot{\sigma(H)} \land  \alpha_f \in \denot{\sigma(F)} \ptag{assumption} \label{found:2-6}
\\& \msteptr{[]}{(\emptyset, e_h)}{\alpha_h}{(\beta, ())} \ptag{assumption} \label{found:2-7}
\\& \msteptr{\alpha_h}{(\beta, \zobs{\overline{x}}{op}{\sigma(e)}))}{\alpha}{(\beta', ())} \ptag{assumption} \label{found:2-8}
\\& \msteptr{\alpha_h \listconcat \alpha }{(\beta', e_f)}{\alpha_f}{(\emptyset, ())} \ptag{assumption} \label{found:2-9}
\\& \exists \alpha'.\exists \overline{c_i}. \alpha = \zevent{op}{\overline{c_i}} \cons \alpha' \land
\alpha_h \models \zevent{op}{\overline{c_i}} \Downarrow \beta_{\eff{op}} \land \notag
\\&\quad\msteptr{\alpha_h \listconcat [\zevent{op}{\overline{c_i}}]}{(\beta \cup \beta_{\eff{op}}, \sigma(e\msubst{x_i}{c_i}))}{\alpha'}{(\beta', ())} \ptag{\textsc{StObs} and \ref{found:2-8}} \label{found:2-10}
\\& [\zevent{op}{\overline{c_i}}] \in \denot{\msgB{op}{\overline{y}}{\sigma(\phi) \land \overline{y = c_i}}} \ptag{lemma~\ref{lemma:lastA-denot}, and $\overline{y} \cap \Code{DOM}(\Gamma) = \emptyset$} \label{found:1-11}
\\& \sigma([\zevent{op}{\overline{x}}])\msubst{x}{c} \in \sigma(\denot{\msgB{op}{\overline{y}}{\phi \land \overline{y = x}}})\msubst{x}{c} \ptag{lift a new substitution $\msubst{x}{c}$ from \ref{found:1-11}} \label{found:1-12}
\\& \sigma(\alpha_h \listconcat [\zevent{op}{\overline{x}}])\msubst{x}{c} \in \sigma(\denot{H \seqA \lastA\msg{op}{\phi \land \overline{y = x}})})\msubst{x}{c} \ptag{lemma~\ref{lemma:seqA-denot}, \ref{found:1-6}, and \ref{found:1-12}} \label{found:2-13}
\end{alignat}}\noindent
Now, we can apply hypothesis~\ref{found:2-5} with
{\footnotesize\begin{align*}
    &\sigma \mapsto \sigma\msubst{x}{c} \quad \alpha_h \mapsto \alpha_h \listconcat [\zevent{op}{\overline{x}}] \quad \alpha_f \mapsto \alpha_f \quad \alpha \mapsto \alpha' \quad \beta \mapsto \beta \cup \beta_{\eff{op}} \quad
    e_h \mapsto e_h; \zobs{\overline{x}}{op}{\sigma(e)} \quad e_f \mapsto e_f
\end{align*}}
Then we have
{\footnotesize\begin{alignat}{2}
&\alpha' \in\denot{\sigma(A\msubst{x}{c})} \ptag{hypothesis~\ref{found:2-5} with \ref{found:2-6}, \ref{found:2-7}, \ref{found:2-9}, \ref{found:2-13}} \label{found:2-14}
\\&\alpha' \in\denot{\sigma(A)} \ptag{$A$ is well formed under context $\Gamma$ and \ref{found:2-14}} \label{found:1-15}
\\&[\zevent{op}{\overline{c_i}}] \listconcat \alpha' \in\denot{\sigma(\lastA\msg{op}{\phi\msubst{x}{v}} \seqA A)}\ptag{hypothesis~\ref{found:1-15}} \label{found:1-16}
\end{alignat}}\noindent
that is sufficient to prove subgoal of this case.

\item:
{\footnotesize
\mprooftr{35mm}
{}
{TRet}
{$\Gamma; \Delta; \emptyset \vdash () : \rg{H}{\epsTL}{F}$}
\ \\ \ \\}
This rule assume that $\Theta \equiv \emptyset, e \equiv (), \tau \equiv \rg{H}{\epsTL}{F}$, thus we need to prove
{\footnotesize\begin{align*}
    \forall \sigma, \sigma \in \denot{\Gamma} \impl \sigma(()) \in \denot{\sigma(\rg{H}{\epsTL}{F})}
\end{align*}}\noindent
that is, prove the term $()$ is in the denotation of a \PAT in from $\rg{H}{\epsTL}{F}$. According to the definition of \PAT denotation, for all$\alpha_h\;\alpha\;\alpha_f\;\beta\;\beta'\;e_h\;e_f$, where $\alpha_h \in \denot{H} \land \alpha_f \in \denot{F}$, we need to show
{\footnotesize\begin{align*}
&\msteptr{[]}{(\emptyset, e_h)}{\alpha_h}{(\beta, ())}
\land \msteptr{\alpha_h}{(\beta, ())}{\alpha}{(\beta', ())}
\land \msteptr{\alpha_h \listconcat \alpha }{(\beta', e_f)}{\alpha_f}{(\emptyset, ())} \impl \alpha \in \denot{\epsTL}
\end{align*}}\noindent
Since there is no small-step reduction rule for the term $()$, thus the relation $\msteptr{\alpha_h}{(\beta, ())}{\alpha}{(\beta', ())}$ is derived from reflexivity case of multi-step reduction. Thus, $\alpha$ is empty trace $[]$, which included by the denotation of $\epsTL$ (\autoref{lemma:espTL}). Then the proof immediate holds in this case.

\item:
{\footnotesize
\mprooftr{21mm}
{
$\Gamma; \Delta; \Theta \vdash e_1 : \tau$ \quad $\Gamma; \Delta; \Theta \vdash e_2 : \tau$
 }
{TChoice}
{$\Gamma; \Delta; \Theta \vdash e_1 \oplus e_2 : \tau$}
\ \\ \ \\}
This rule assume that $e \equiv e_1 \oplus e_2$, thus we need to prove
{\footnotesize\begin{align*}
    \forall \sigma, \sigma \in \denot{\Gamma} \impl \sigma(e_1 \oplus e_2) \in \denot{\sigma(\tau)}
\end{align*}}\noindent
From the inductive hypothesis of this case, we know
{\footnotesize\begin{align*}
    &\forall \sigma, \sigma \in \denot{\Gamma} \impl \sigma(e_1) \in \denot{\sigma(\tau)}
    \\&\forall \sigma, \sigma \in \denot{\Gamma} \impl \sigma(e_2) \in \denot{\sigma(\tau)}
\end{align*}}\noindent
Then the Lemma~\ref{lemma:choice-denot} is sufficient to prove the subgoal of this case.

\item:
{\footnotesize
\mprooftr{50mm}
{$\Gamma, z\hasoty{unit}{\phi};\Delta;\Theta \vdash e : \tau$ \quad $z$ is fresh}
{TAssume}
{$\Gamma;\Delta;\Theta  \vdash \zassume{\phi}{e} : \tau  $}
\ \\ \ \\}
This rule assume that $e \equiv \zassume{\phi}{e}$, thus we need to prove
{\footnotesize\begin{align*}
    \forall \sigma, \sigma \in \denot{\Gamma} \impl \sigma(\zassume{\phi}{e}) \in \denot{\sigma(\tau)}
\end{align*}}\noindent
From the inductive hypothesis of this case, we know
{\footnotesize\begin{align*}
    &\forall \sigma, \sigma \in \denot{\Gamma, z\hasoty{unit}{\phi}} \impl \sigma(e) \in \denot{\sigma(\tau)}
\end{align*}}\noindent
Since $z$ is a fresh variable, then we have
{\footnotesize\begin{align*}
    &\forall \sigma, \sigma \in \denot{\Gamma, z\hasoty{unit}{\phi}} \impl \exists \sigma'. \sigma'[\overline{z \mapsto ()}] = \sigma. \sigma'(e) \in \denot{\sigma'(\tau)}
\end{align*}}\noindent
Moreover, according to the definition of type context denotation,
{\footnotesize\begin{align*}
    &\forall \sigma, \sigma \in \denot{\Gamma} \land \sigma(\phi) \iff \sigma[\overline{z \mapsto ()}] \in \denot{\Gamma, z\hasoty{unit}{\phi}}
\end{align*}}\noindent
Then it is safe to apply Lemma~\ref{lemma:pure-denot} with $\sigma$ as substitution in $\denot{\Gamma}$ and make $\sigma(\phi)$ holds, and $e \mapsto \sigma(\zassume{\phi}{e}), e' \mapsto \sigma(e), \tau \mapsto \sigma(\tau)$.
Now, we need to show $\zassume{\phi}{e}$ can reduced into $e$ without add new effect, which is can be proved by \textsc{StAssume} and $\sigma(\phi)$. Then the proof immediate holds in this case.

\item:
{\footnotesize
\mprooftr{45mm}
{$\Gamma;\Delta;\Theta \vdash e : \tau$ \quad
$\Gamma \vdash () : \nuot{unit}{\phi}$}
{TAssert}
{$\Gamma;\Delta;\Theta \vdash \zassert{\phi}{e} : \tau  $}
\ \\ \ \\}
This rule assume that $e \equiv \zassert{\phi}{e}$, thus we need to prove
{\footnotesize\begin{align*}
    \forall \sigma, \sigma \in \denot{\Gamma} \impl \sigma(\zassert{\phi}{e}) \in \denot{\sigma(\tau)}
\end{align*}}\noindent
From the assumption and inductive hypothesis of this case, we know
{\footnotesize\begin{align*}
    &\forall \sigma, \sigma \in \denot{\Gamma} \impl \sigma(e) \in \denot{\sigma(\tau)} \land \sigma(\phi)
\end{align*}}\noindent
Then it is safe to apply Lemma~\ref{lemma:pure-denot} with $e \mapsto \sigma(\zassert{\phi}{e}), e' \mapsto \sigma(e), \tau \mapsto \sigma(\tau)$.
Now, we need to show $\zassert{\phi}{e}$ can reduced into $e$ without add new effect, which is can be proved by \textsc{StAssert} and $\sigma(\phi)$. Then the proof immediate holds in this case.

\item:
{\footnotesize
\mprooftr{58mm}
{
$\Gamma \vdash \primop : t$ \quad $\Gamma \vdash t <: \overline{y{:}t}\sarr t_x$
$\forall i. \Gamma \vdash v_i : t_i$ \quad
$\Gamma, x{:}t_x\msubst{y}{v};\Delta;\Theta \vdash e : \tau$
}
{TOpApp}
{$\Gamma;\Delta;\Theta \vdash \zlet{x{:}b}{\primop\ \overline{v}}{e} : \tau$}
\ \\ \ \\}
This rule assume that $e \equiv \zlet{x{:}b}{\primop\ \overline{v}}{e}$, thus we need to prove
{\footnotesize\begin{align*}
    \forall \sigma, \sigma \in \denot{\Gamma} \impl \sigma(\zlet{x{:}b}{\primop\ \overline{v}}{e}) \in \denot{\sigma(\tau)}
\end{align*}}\noindent
From the assumption and inductive hypothesis of this case, we know
{\footnotesize\begin{alignat}{2}
\setcounter{equation}{0}
&\Delta(\eff{op}) = t \ptag{assumption} \label{found:3-1} \\
&\Gamma \vdash t <: \overline{y{:}t}\sarr t_x \ptag{assumption} \label{found:3-2} \\
&\forall i. \Gamma \vdash v_i : t_i \ptag{assumption} \label{found:3-3} \\
&\Gamma, x{:}t_x\msubst{y}{v};\Delta;\Theta \vdash e : \tau \ptag{assumption} \label{found:3-4} \\
&\forall v_x. \Gamma \vdash v_x : t_x\msubst{y}{v} \impl \Gamma;\Delta;\Theta \vdash e[x\mapsto v_x] : \tau[x\mapsto v_x] \ptag{Lemma~\ref{lemma:subt} and \ref{found:3-4}} \label{found:3-5} \\
&\forall \sigma \in \denot{\Gamma}. \forall v_x \in \denot{\sigma(t_x\msubst{y}{v})}. \sigma(e[x\mapsto v_x]) \in \denot{\sigma(\tau[x\mapsto v_x])} \ptag{induction hypothesis and \ref{found:3-5}} \label{found:3-6}
\\& \forall i. \forall \sigma\in\denot{\Gamma}. \sigma(v_i) \in\denot{\sigma(t_i)} \ptag{Lemma~\ref{theorem:pure-fundamental-ap} and \ref{found:3-3}} \label{found:3-7}
\\& \forall \sigma\in\denot{\Gamma}. \forall c_x. \primop(\overline{\sigma(v)}) \Downarrow c_x \impl c_x \in \denot{t_x\msubst{y}{v}} \ptag{Lemma~\ref{lemma:built-in-typing-sub-pure}, \ref{found:3-1}, \ref{found:3-2}, and \ref{found:3-6}} \label{found:3-8}
\\& \forall \sigma\in\denot{\Gamma}. \forall c_x. \primop(\overline{\sigma(v)}) \Downarrow c_x \impl \sigma(e[x\mapsto c_x]) \in \denot{\sigma(\tau[x\mapsto c_x])} \ptag{\ref{found:3-6} and \ref{found:3-8}} \label{found:3-9}
\\& \forall \sigma\in\denot{\Gamma}. \forall c_x. \primop(\overline{\sigma(v)}) \Downarrow c_x \impl \sigma(e[x\mapsto c_x]) \in \denot{\sigma(\tau)} \ptag{\ref{found:3-9} and $\tau$ is well-formed under $\Gamma$} \label{found:3-10}
\end{alignat}}\noindent
Then it is safe to apply Lemma~\ref{lemma:pure-denot} with $e \mapsto \sigma(\zlet{x{:}b}{\primop\ \overline{v}}{e}), e' \mapsto \sigma(e[x\mapsto c_x]), \tau \mapsto \sigma(\tau)$.
Now, we need to show $\zlet{x{:}b}{\primop\ \overline{v}}{e}$ can reduced into $e[x\mapsto c_x]$ without add new effect, which is can be proved by \textsc{StOp} and the assumption $\primop(\overline{\sigma(v)}) \Downarrow c_x$. Then the proof immediate holds in this case.

\item:
{\footnotesize
\mprooftr{38mm}
{$\Gamma;\Delta;\Theta \vdash e : \tau$ \quad
$\Gamma \vdash \tau <: \tau'$}
{TSub}
{$\Gamma;\Delta;\Theta \vdash e : \tau'$}
\ \\ \ \\}
The case can be directly proved by Lemma~\ref{lemma:subtyping}.
\end{enumerate}
\end{proof}

